\documentclass[10pt,a4paper]{amsart}
\usepackage[margin=19mm]{geometry}
\usepackage{amsmath,amssymb,mathtools}
\usepackage[scr=rsfs,cal=euler]{mathalfa}
\usepackage{xfrac}
\usepackage[unicode,hidelinks,bookmarks=false]{hyperref}
\newcommand{\ie}{i.e.\ }
\newcommand{\cf}{cf.\ }
\newcommand{\resp}{respectively\ }
\newcommand{\Subsection}{\subsection}
\providecommand{\nobreakdash}{\nobreak\hbox{-}}
\setlength{\emergencystretch}{2em}
\multlinegap0pt
\usepackage{enumerate}
\usepackage{amssymb}
\newtheorem{prop}{Proposition}
\title{Perturbations of Cauchy differences}
\author{Eszter Gselmann, Tomasz Małolepszy and Janusz Matkowski}
\date{Source: arXiv:2603.19242v1 (CC BY 4.0)}
\begin{document}
\maketitle

\noindent\emph{Source and licence.} Perturbations of Cauchy differences. Version arXiv:2603.19242v1 (CC BY 4.0), distributed by arXiv under the Creative Commons Attribution 4.0 International licence, http://creativecommons.org/licenses/by/4.0/, as recorded in the arXiv licence metadata for this exact version and shown on the abstract page https://arxiv.org/abs/2603.19242v1. Full licence notice: \texttt{licenses/arxiv-2603.19242-CC-BY-4.0.txt}.\par\smallskip

\noindent\textit{Editorial scope.} Counted unit: the article's prop prop3 (source0.tex lines 521--545). The article's complete original proof of it is delivered with it (source0.tex lines 547--612, one \texttt{proof} environment of 66 lines). No mathematics is written, repaired or re-derived: the statement and the proof are the article's own lines, in the article's own notation, and no retained line is substituted or re-flowed.  Retained uncounted as the source-local context the proof needs: lines 285--299.  Nothing was omitted from any retained span, so the package carries no omission record.  No retained line contains a citation, so the wrapper carries no bibliography and no bibliography entry.  No retained line contains a figure environment, an image-inclusion command or drawing code, so no image is delivered and the package declares no asset dependency.  Editorial formatting only, in addition: an AMS \texttt{amsart} preamble that restates the article's own declarations and macros used by the retained lines, namely the environments \texttt{prop} on separate counters, together with the standard packages those lines need.  The retained environments print in this document as Proposition 1 in order of appearance, whereas the article prints the same items with the numbering of its own full document.  The complete native source of the article as distributed in the arXiv e-print for this exact version is bundled unchanged as \texttt{sources/196763/source0.tex}.
\begin{prop}\label{prop_biadditive}
 Let $(S, +)$ be a semigroup, $(\mathbb{F},+, \cdot, )$ be a field, and $B\colon S\times S\to \mathbb{F}$ be a symmetric and biadditive function. If a function $f\colon S\to \mathbb{F}$ satisfies equation 
 \begin{equation}\label{biadditive}
  f(x+y)-f(x)-f(y)= B(x, y)
 \end{equation}
for all $x, y\in S$, then exists an additive function $a\colon S\to \mathbb{F}$ i.e., a function for which 
\[
 a(x+y)= a(x)+a(y)
\]
holds for all $x, y\in S$ such that 
\[
 f(x)= \frac{1}{2}B(x, x)+a(x)
\]
for all $x\in S$. 
\end{prop}

\begin{prop}\label{prop3}
 Let $(\mathbb{F}, +, \cdot)$ be a field. If the functions $f, \alpha \colon \mathbb{F}\to \mathbb{F}$ satisfy 
 \begin{equation}\label{add_product}
  f(x+y)-f(x)-f(y)= \alpha(xy)
 \end{equation}
for all $x, y\in \mathbb{F}$, then there exist a constant $\gamma\in \mathbb{F}$, and there are additive functions $a, A\colon \mathbb{F}\to \mathbb{F}$, i.e., functions for which 
\[
 a(x+y)= a(x)+a(y) 
 \quad 
 \text{and}
 \quad 
 A(x+y)= A(x)+A(y) 
 \qquad 
 \left(x, y\in \mathbb{F}\right)
\]
such that 
\[
 \alpha(x)= a(x)+\gamma 
 \quad 
 \text{and}
 \quad 
 f(x)= \frac{1}{2}a(x^{2})+A(x)+\gamma
\]
holds for all $x\in \mathbb{F}$. 
\end{prop}

\begin{proof}
 Let us consider the function $\mathscr{C}_{f}\colon \mathbb{F}\times \mathbb{F}\to \mathbb{F}$ defined by 
 \[
\mathscr{C}_{f}(x, y)= f(x+y)-f(x)-f(y)  
\qquad 
\left(x, y\in \mathbb{F}\right). 
 \]
 As $\mathscr{C}_{f}$ is a Cauchy difference generated by the function $f$, it fulfills the cocycle equation on $\mathbb{F}\times \mathbb{F}$, i.e., we have 
 \[
  \mathscr{C}_{f}(x, y)+\mathscr{C}_{f}(x+y, z)= \mathscr{C}_{f}(x, y+z)+\mathscr{C}_{f}(y, z) 
  \qquad 
  \left(x, y, z\in \mathbb{F}\right). 
 \]
 Further, due to our equation, we have $\mathscr{C}_{f}(x, y)= \alpha(xy)$ for all $x, y\in \mathbb{F}$. Therefore 
 \[
  \alpha(xy)+\alpha((x+y)z)= \alpha(x(y+z))+\alpha(yz) 
  \qquad 
  \left(x, y, z\in \mathbb{F}\right), 
 \]
or after some rearrangement 
\[
 \alpha(xy)+\alpha(xz+yz)= \alpha(xy+xz)+\alpha(yz). 
\]
In other words, we have 
\[
 \alpha(u)+\alpha(v+w)= \alpha(u+v)+\alpha(w) 
 \qquad 
 \left(u, v, w\in \mathbb{F}\right). 
\]
Therefore, there exists a constant $\gamma\in \mathbb{F}$ and an additive function $a\colon \mathbb{F}\to \mathbb{F}$ such that 
\[
 \alpha(x)= a(x)+\gamma 
 \qquad 
 \left(x\in \mathbb{F}\right). 
\]
Define the function $B\colon \mathbb{F}\times \mathbb{F}\to \mathbb{F}$ through 
\[
 B(x, y)= a (xy) 
 \qquad 
 \left(x, y\in \mathbb{F}\right). 
\]
Due to the additivity of the function $a$, the mapping $B$ is symmetric and biadditive. Consider now the function $\tilde{f}$ defined on $\mathbb{F}$ by 
\[
 \tilde{f}(x)= f(x)-\gamma 
 \qquad 
 \left(x\in \mathbb{F}\right)
\]
to deduce that we have 
\[
 \tilde{f}(x+y)-\tilde{f}(x)-\tilde{f}(y)= B(x, y) 
 \qquad 
 \left(x, y\in \mathbb{F}\right). 
\]
Observe that Proposition \ref{prop_biadditive} can be applied with the choice $S=\mathbb{F}$ to obtain that there exists an additive function $A\colon \mathbb{F}\to \mathbb{F}$ such that 
\[
 \tilde{f}(x)= \frac{1}{2}B(x, x)+A(x)=
 \frac{1}{2} a(x^{2})+A(x) 
 \qquad 
 \left(x\in \mathbb{F}\right).
\]
From this however we conclude that 
\[
 f(x)= \tilde{f}(x)+\gamma= \frac{1}{2} a(x^{2})+A(x)+\gamma
\]
for all $x\in \mathbb{F}$. 
\end{proof}

\end{document}
